\chapter{Conclusiones y trabajo futuro}
\label{chap:conclusiones}

Este Trabajo Fin de Máster ha tratado la creación de lenguajes de dominio
específico con sintaxis basada en bloques. La motivación parte de un problema
práctico: Blockly facilita la creación de editores visuales, pero cada dominio
todavía requiere definir a mano bloques, campos, toolbox, conexiones,
validaciones y generadores. El trabajo propone una alternativa basada en
ingeniería de lenguajes y modelos: describir el dominio a alto nivel y generar
automáticamente el editor Blockly.

El capítulo se organiza en torno a los resultados y las perspectivas de la
propuesta. En primer lugar, se analiza la viabilidad del enfoque y se resumen
sus principales aportaciones. A continuación, se revisan las conclusiones
obtenidas para cada objetivo y se presentan las limitaciones identificadas.
Después se proponen posibles líneas de trabajo futuro y, finalmente, se ofrece
un cierre global del trabajo realizado.

\section{Viabilidad del enfoque}
\label{sec:viabilidad-enfoque}

El trabajo desarrollado aporta evidencia de que es viable diseñar e implementar
un entorno para la creación de lenguajes de dominio específico con sintaxis
basada en bloques. Model2Blockly acepta metamodelos Ecore enriquecidos con
anotaciones y modelos textuales \texttt{.m2b} definidos con Xtext. Ambas rutas
se transforman en un modelo intermedio común, y desde él se generan editores
Blockly ejecutables. La herramienta se materializa además como plugin Eclipse
instalable mediante un update site p2, por lo que el enfoque no queda limitado a
un prototipo ejecutado manualmente desde el código fuente.

La solución aprovecha información que ya existe en una descripción de dominio.
Las clases se convierten en tipos de bloque, los atributos en campos, las
contenciones en entradas estructurales, las referencias en desplegables
dinámicos y las cardinalidades en comprobaciones. Las anotaciones añaden
metadatos de presentación. Este proceso reduce la programación manual en
JavaScript y sigue la idea central de la ingeniería dirigida por modelos:
utilizar modelos para especificar y generar software
\cite{brambilla2017modelDriven,hutchinson2014industry}.

El trabajo también refuerza la utilidad de separar lo que el dominio significa
de la forma en que se muestra al usuario. El dominio se describe mediante Ecore
o mediante el DSL textual \texttt{.m2b}; Blockly actúa como la representación
visual generada para manipular ese dominio. Esta separación es coherente con la
forma en que los lenguajes de dominio específico organizan sus conceptos y sus
notaciones \cite{sal2024dsl}.

\section{Principales aportaciones}
\label{sec:aportaciones}

La primera aportación es una arquitectura generativa que conecta Ecore,
\texttt{.m2b}, \texttt{EditorSpec} y Blockly. La arquitectura introduce un
modelo intermedio, en lugar de acoplar directamente cada entrada a una salida
concreta. Las dos rutas se normalizan hacia el mismo contrato y no requieren
generadores separados.

La segunda aportación es el modelo intermedio común. Este
modelo recoge los elementos necesarios para construir un editor visual: bloques,
campos, categorías, entradas de valor, entradas de sentencias, referencias,
validaciones, opciones de workspace y metadatos de código. También define el
contrato entre adaptadores y generadores, se serializa como XMI y se recarga
antes de generar la salida final.

La tercera aportación es la ruta Ecore. El adaptador desarrollado permite partir
de un metamodelo y derivar una especificación de editor. Esta ruta cubre
la parte metamodelo-céntrica del sistema: usar un metamodelo, enriquecido cuando
sea necesario con anotaciones, para generar código basado en Blockly.

La cuarta aportación es la ruta textual \texttt{.m2b}. Esta ruta facilita definir
lenguajes de bloques de forma compacta, con una sintaxis cercana a los conceptos
que después se convertirán en bloques. También incorpora validación estática
mediante Xtext, lo que permite detectar errores antes de generar el editor cuando
se usa esta forma de autoría.

La quinta aportación es la generación de editores ejecutables. La salida
HTML/JavaScript produce definiciones de bloques, toolbox, generadores,
validaciones, páginas HTML autocontenidas, exportación JSON/XMI, informes de
generación, modelos de ejemplo y una vista Blockly de las reglas de validación.

La sexta aportación es la integración práctica de la herramienta. El plugin
Eclipse ofrece comandos para generar desde ficheros \texttt{.m2b} y
\texttt{.ecore}, y para aplicar de vuelta a la fuente cambios soportados en las
reglas visuales de validación. El update site y las demos publicadas facilitan la
revisión externa del resultado.

Finalmente, el trabajo aporta una evaluación basada en un ejemplo conductor y
pruebas automáticas. AppMaker se describe como metamodelo Ecore anotado y como
fichero textual \texttt{.m2b}, lo que permite comparar las dos formas de autoría.
Las 144 pruebas JUnit organizadas por capa aportan evidencia adicional sobre el
comportamiento del modelo intermedio, los adaptadores y los generadores.

\section{Conclusiones por objetivos}
\label{sec:conclusiones-objetivos}

Los objetivos específicos definidos al inicio del trabajo se han abordado de la
siguiente forma:

\begin{itemize}
  \item \textbf{OE1.} Se analizaron los elementos necesarios para un editor de
  bloques y se representaron explícitamente en el modelo intermedio: tipos de
  bloque, campos, conexiones, categorías, referencias, validaciones y
  generación.
  \item \textbf{OE2.} Se implementó una ruta Ecore que transforma clases,
  atributos, referencias, herencia, cardinalidades y anotaciones en una
  especificación Blockly.
  \item \textbf{OE3.} Se definió una ruta textual \texttt{.m2b} con Xtext para
  describir lenguajes de bloques sin editar directamente ficheros Ecore en
  ejemplos o configuraciones compactas.
  \item \textbf{OE4.} Se diseñó e implementó un modelo intermedio común como
  representación común independiente de la fuente de entrada.
  \item \textbf{OE5.} Se generaron editores HTML/JavaScript ejecutables con
  bloques, toolbox, validaciones, guardado, carga, exportación JSON/XMI e
  informes de generación.
  \item \textbf{OE6.} Se incorporaron comprobaciones, referencias dinámicas,
  personalización visual, metadatos de interfaz, plantillas de código y una
  vista Blockly editable para reglas de validación.
  \item \textbf{OE7.} Se evaluó la solución con AppMaker por la ruta Ecore y por
  la ruta textual \texttt{.m2b}, junto con pruebas automáticas sobre las capas
  principales del sistema.
\end{itemize}

Por lo tanto, el enfoque es técnicamente viable dentro del alcance definido.
Un dominio descrito con Ecore anotado o con \texttt{.m2b} puede transformarse en
un editor de bloques mediante una cadena generativa reproducible. La ruta
textual muestra además que esa información puede escribirse de forma más
compacta y normalizarse hacia el mismo modelo intermedio.

\section{Limitaciones}
\label{sec:conclusiones-limitaciones}

Aunque los resultados respaldan el enfoque principal, el trabajo tiene
limitaciones. La primera es que la evaluación usa un ejemplo diseñado dentro del
proyecto. AppMaker cubre estructuras variadas, pero no equivale a una validación
industrial con muchos metamodelos externos. Por ello, el caso estudiado apoya la
viabilidad de la herramienta, pero no prueba su generalidad de forma exhaustiva.

La segunda limitación afecta a las validaciones. Model2Blockly genera advertencias
para mínimos y máximos de bloques, campos obligatorios, referencias obligatorias,
unicidad, reglas simples de orden, expresiones declaradas en anotaciones y un
subconjunto sencillo de invariantes OCL. Este soporte cubre restricciones
frecuentes, pero no equivale a un lenguaje completo de restricciones como OCL
\cite{omgOcl}. No se soportan invariantes arbitrarios, cuantificadores,
navegación compleja ni análisis semántico profundo. Además, las comprobaciones
generadas producen advertencias, no una validación bloqueante. La edición visual
de reglas tampoco es una transformación bidireccional general: solo se aplican de
vuelta a la fuente los cambios que tienen una correspondencia clara.

La tercera limitación se refiere a las referencias. Las referencias no
contenidas se representan mediante desplegables dinámicos, lo que permite
seleccionar elementos existentes del workspace. Además, el runtime HTML
sincroniza referencias opuestas no contenidas cuando ambos extremos son
editables. La limitación actual está en los casos más complejos: referencias
opuestas con semántica de contenedor y visualización avanzada de relaciones
entre elementos.

La cuarta limitación está relacionada con la generación de código. La solución
incluye plantillas de código y una representación textual por defecto, pero no
incorpora todavía formateadores específicos, análisis de tipos avanzado,
gestión de imports, generación multiarchivo ni integración directa con
herramientas externas de construcción o ejecución.

La quinta limitación es de evaluación de usuario. Este TFM no mide si los
editores generados reducen efectivamente el esfuerzo de aprendizaje o mejoran la
experiencia de usuarios no técnicos. La literatura sobre programación por
bloques muestra que la modalidad de interacción influye en las prácticas de los
usuarios noveles \cite{weintrop2018modalities}, pero validar empíricamente los
editores generados requeriría un estudio específico.

\section{Trabajo futuro}
\label{sec:trabajo-futuro}

A partir de estas limitaciones se identifican varias líneas de trabajo futuro.

En primer lugar, sería útil ampliar el soporte de validación. Una opción sería
integrar un lenguaje de restricciones más expresivo o un mecanismo de reglas
declarativas. Así se podrían describir condiciones más complejas que los mínimos,
máximos y reglas de orden inmediato. Esta línea también permitiría ampliar la
traducción entre reglas visuales editadas y su representación en
\texttt{.m2b} o en anotaciones Ecore.

En segundo lugar, la gestión de referencias podría evolucionar hacia un soporte
más completo de referencias bidireccionales. Esto implicaría extender la
sincronización de referencias opuestas en el runtime HTML, cubrir casos de
contenedor, detectar referencias que apuntan a elementos ya eliminados y ofrecer
visualizaciones más claras de las relaciones entre elementos del workspace.

En tercer lugar, se podría profundizar en la importación y exportación de
modelos. Actualmente la herramienta se centra en generar editores y exportar el
contenido del workspace. Una línea futura consistiría en importar modelos XMI
existentes y reconstruir automáticamente el workspace Blockly correspondiente.
Esto acercaría más la herramienta a flujos completos de ingeniería de modelos.

En cuarto lugar, la generación de código podría enriquecerse. Las plantillas
actuales son suficientes para demostrar una ruta genérica de exportación, pero
podrían ampliarse con validación de plantillas, tipado, formateadores,
generación multiarchivo y objetivos específicos para lenguajes o plataformas
concretas.

En quinto lugar, podrían explorarse interfaces web más avanzadas sobre el mismo
modelo intermedio. Un futuro objetivo de generación podría incorporar un
inspector estructurado, temas visuales por dominio, componentes personalizados y
modos de edición más adaptados a usuarios finales.

En sexto lugar, sería conveniente realizar una evaluación empírica con usuarios.
Un estudio futuro podría comparar el esfuerzo necesario para construir un editor
Blockly manualmente frente al uso de Model2Blockly, o analizar cómo interactúan
usuarios no técnicos con editores generados para distintos dominios. Esta
evaluación permitiría complementar la evidencia técnica de este TFM con datos
de uso.

Por último, convendría robustecer el ciclo de publicación del plugin. La versión
actual ya dispone de update site p2 y verificaciones de empaquetado, pero una
línea futura sería automatizar más la publicación, el versionado y las pruebas de
instalación sobre instalaciones limpias de Eclipse.

\section{Cierre}
\label{sec:cierre}

Model2Blockly aporta evidencia de que es posible conectar técnicas de ingeniería de
lenguajes con entornos de programación basada en bloques. El resultado no es un
único editor visual, sino una infraestructura que transforma descripciones de
dominio en editores Blockly ejecutables. El conocimiento del dominio se mantiene
en modelos o DSLs de alto nivel, y la sintaxis visual se genera a partir de ellos.

El trabajo muestra que la combinación de Ecore, Xtext, un modelo intermedio EMF,
generación HTML/JavaScript y Blockly es viable para construir editores de
lenguajes de dominio específico con sintaxis basada en bloques. Dentro del
alcance definido, esta combinación es una alternativa razonable a la programación
manual de editores Blockly y deja una base extensible para trabajos futuros.
